% flags-experiments-observability.tex — feature flags (release / experiment / ops /
% permission, remote config, kill switch, local overrides, evaluation at launch +
% caching, flag debt), A/B testing (randomisation unit, hashed bucketing, exposure
% logging, SRM, sample size / power / significance, peeking, guardrails), analytics
% event design (schema, naming, PII, batching + offline queue), mobile observability
% (crash-free users vs sessions, symbolication, hangs/ANR, traces, breadcrumbs),
% staged rollout (App Store phased release) and release-health gates.
% Sources (checked 2026-09-25) — phased release; the rest are published works: Pete Hodgson, "Feature Toggles (aka Feature Flags)", martinfowler.com
% (the four toggle categories); Kohavi, Tang, Xu, "Trustworthy Online Controlled
% Experiments" (SRM, peeking, guardrails, triggering); Lehr's rule n = 16 s^2/d^2
% (alpha 0.05 two-sided, power 0.8); App Store Connect Help "Release a version
% update in phases" (checked 2026-09-25: 1/2/5/10/20/50/100 % over 7 days,
% automatic updates only, manual download any time, pause up to 30 days total,
% release to all any time); SE-0206 (Hasher seeded per process). Vendors named
% only as examples, no vendor specifics.
% Not repeated: dSYM upload + crash report anatomy (security-build/crashes-symbolication),
% TestFlight / fastlane pipeline (security-build/signing-and-cicd), MetricKit / Organizer
% (debugging/logging-signposts-power).
% @labels: area=engineering kind=process platform=general topic=system-design,data,debugging discussion=129
% @tags: feature-flags, kill-switch, ab-testing, hashed-bucketing, exposure-logging, sample-ratio-mismatch, sample-size, peeking, guardrail-metrics, analytics-events, crash-free-users, phased-release
\documentclass[8pt]{extarticle}
\usepackage{printup-sheet}
\usepackage{array}

\lstdefinelanguage{SwiftSheet}{
  morekeywords={import,struct,final,class,var,let,func,return,private,init,try,
    async,await,self,true,false,String,Int,UInt32,Data},
  sensitive=true, morecomment=[l]{//}, morestring=[b]",
  literate={??}{{\hbox{?}\hbox{?}}}2 {->}{{\hbox{-}\hbox{>}}}2 {<<}{{\hbox{<}\hbox{<}}}2
           {==}{{\hbox{=}\hbox{=}}}2 {!=}{{\hbox{!}\hbox{=}}}2}
\lstset{basicstyle=\ttfamily\footnotesize, aboveskip=2pt, belowskip=2pt}

\newcommand\ct[1]{\texttt{#1}}

\tikzset{
  lp/.style={box, font=\scriptsize, text width=25mm, align=center, inner sep=2pt, minimum height=11mm},
  lbl/.style={font=\tiny, text=black!75, inner sep=1pt, align=center},
  hd/.style={font=\bfseries\small, text=sheetBlue, anchor=west},
}

\begin{document}

\sheettitle{Feature flags, experiments \& mobile observability}{engineering · folio}

\oneliner{Ship code \textbf{dark} behind a \textbf{flag}, turn it on for a
\textbf{hashed, random slice} of users, \textbf{log who was exposed}, compare their
metrics with control while \textbf{guardrails} (crashes, hangs, revenue) watch for
harm, then ramp, hold or \textbf{kill} — without a new binary. On mobile this matters
twice: a shipped build lives for months and the store rollout cannot be rolled back.}

\vspace{2pt}
\noindent\begin{tikzpicture}[sheet]
  \node[hd] at (-0.1,3.05) {The loop: flag + experiment + telemetry};
  \node[lp] (a) at (1.25,2.3) {\textbf{1 ship dark}\\code + flag, default \emph{off}, compiled-in default};
  \node[lp, draw=sheetBrown, fill=sheetBrown!7] (b) at (4.25,2.3) {\textbf{2 remote config}\\rules, \% ramp, arms: \ct{hash(exp:user)}};
  \node[lp, draw=sheetGreen!70!black, fill=sheetGreen!8] (c) at (7.25,2.3) {\textbf{3 app evaluates}\\cached at launch, stable for the session};
  \node[lp, draw=sheetOrange, fill=sheetOrange!8] (d) at (7.25,0.25) {\textbf{4 telemetry}\\exposure + events, crashes, hangs, traces};
  \node[lp, draw=sheetOrange, fill=sheetOrange!8] (e) at (4.25,0.25) {\textbf{5 analyse}\\SRM check, primary metric, guardrails, CI};
  \node[lp, draw=sheetRed, fill=sheetRed!6] (f) at (1.25,0.25) {\textbf{6 decide}\\ramp · hold · \textbf{kill switch} · remove flag};
  \draw[flow] (a) -- (b); \draw[flow] (b) -- (c); \draw[flow] (c) -- (d);
  \draw[flow] (d) -- (e); \draw[flow] (e) -- (f);
  \coordinate (x1) at (0.7,0); \coordinate (x2) at (1.9,0); \coordinate (x3) at (3.7,0);
  \draw[flow] (f.north -| x1) -- node[lbl, right, align=left]{delete\\flag} (a.south -| x1);
  \draw[hot] (f.north -| x2) |- node[lbl, above, pos=0.75]{ramp / kill:\\config, no release} (3.7,1.28) -- (b.south -| x3);
  \draw[sheetGrey!50] (9.0,3.2) -- (9.0,-0.2);
  % ── phased release ──
  \node[hd] at (9.1,3.05) {App Store phased release (automatic updates)};
  \foreach \d/\p in {1/1,2/2,3/5,4/10,5/20,6/50,7/100} {
    \pgfmathsetmacro\h{0.02*\p}
    \fill[sheetBlue!55] (9.05+\d*0.82-0.25,0.35) rectangle ++(0.5,\h);
    \node[lbl, anchor=south] at (9.05+\d*0.82,0.35+\h) {\p\,\%};
    \node[lbl, anchor=north] at (9.05+\d*0.82,0.33) {day \d};
  }
  \draw[sheetGrey] (9.3,0.35) -- (15.35,0.35);
  \foreach \d in {1,...,6} \node[font=\tiny, text=sheetGreen!45!black] at (9.05+\d*0.82+0.41,1.1) {\checkmark};
  \node[lbl, anchor=west, align=left, text=sheetGreen!45!black] at (9.15,1.55) {gate before each step:\\crash-free users, hang rate,\\funnel vs previous version};
  \node[lbl, anchor=north west, align=left] at (9.1,-0.05) {pause up to 30 days in total · release to all any time · manual downloads get it at once ·\\\textbf{no rollback}: a bad build needs a new build or a server-side kill switch};
\end{tikzpicture}

\begin{multicols}{2}
\small\setstretch{1.0}\raggedright

\section{Feature flags}
{\footnotesize\setlength{\tabcolsep}{3pt}
\begin{tabular}{@{}>{\raggedright\arraybackslash}p{13mm}>{\raggedright\arraybackslash}p{30mm}>{\raggedright\arraybackslash}p{32mm}@{}}
\toprule
\textbf{kind} & \textbf{purpose} & \textbf{lifetime · decided by}\\ \midrule
release & hide unfinished work, trunk-based dev & days–weeks · per build/env\\
experiment & A/B arms & weeks · per user, hashed\\
ops & kill switch, degrade under load & short or permanent · ops, at runtime\\
permission & premium, beta, internal & long · per user/entitlement\\
\bottomrule
\end{tabular}}
\begin{itemize}
  \item \textbf{Evaluation}: compiled-in defaults (first launch, offline) $\to$
        last-known-good cache on disk $\to$ fetch at launch with a short
        timeout, applied at a \emph{safe point} (next launch / before the screen
        builds) — never flip UI mid-session. Experiments read once per session.
  \item \textbf{Kill switch}: server-side off for a risky path without a
        release; test the \emph{off} path too. Add a minimum-version /
        force-update switch before you need it.
  \item \textbf{Local overrides}: debug menu, launch arguments (\ct{-newCheckout
        YES} lands in \ct{UserDefaults}' argument domain), UI tests set flags
        explicitly — never depend on the live server.
  \item \textbf{Flag debt}: owner + expiry per flag, delete at 100\,\%, lint for
        stale keys. Mobile twist: old app versions still read the key — keep a
        server default until they are gone. $n$ flags = $2^n$ paths.
\end{itemize}

\section{A/B testing}
\begin{itemize}
  \item \textbf{Randomisation unit} = analysis unit: user id (stable across
        devices) > install id (resets on reinstall) > session.
  \item \textbf{Bucketing}: \ct{hash(experimentKey + ":" + userId) mod 10000};
        arm = range. Deterministic, sticky, stateless; the salt makes
        experiments independent; ramping by widening a range keeps users in place.
  \item \textbf{Exposure logging}: log \ct{experiment\_exposure} when the user
        \emph{reaches} the changed code, not at assignment; analyse exposed users
        only (otherwise the effect is diluted).
  \item \textbf{SRM} (sample-ratio mismatch): 50/50 planned, 50.8/49.2 at large
        $n$ $\to$ $\chi^2$ fails $\to$ assignment or logging bug; results void.
  \item \textbf{Sample size} ($\alpha$ = 0.05 two-sided, power 80\,\%):
        $n_{\mathrm{arm}} \approx 16\,\sigma^2/\delta^2$; conversion
        $\sigma^2 = p(1-p)$. 10\,\% $\to$ 11\,\%: $16 \cdot 0.09/0.01^2 \approx
        14\,400$ per arm. Smaller effect $\to$ quadratically more users.
  \item \textbf{Peeking}: checking daily and stopping at the first $p<0.05$
        inflates false positives far above 5\,\%. Fix the horizon (full weeks),
        or use a sequential test built for it.
  \item \textbf{Guardrails}: crash-free users, latency, revenue, uninstalls
        must not regress even when the primary metric wins. Many metrics
        $\to$ some ``win'' by chance.
\end{itemize}

\section{Analytics event design}
\begin{itemize}
  \item \textbf{Name} \ct{object\_action}, past tense, snake\_case:
        \ct{checkout\_started}, \ct{item\_added}. A \textbf{tracking plan}
        (schema registry) types every property; validate in debug and CI.
  \item \textbf{Envelope}: \ct{event\_id} (UUID, server dedupe), client time +
        server receive time (skewed clocks), anonymous/user id, session id, app
        version + build, OS, device, active experiment arms.
  \item \textbf{No PII}: no email, name, free text or precise location;
        pseudonymous ids; honour consent (ATT for cross-app tracking). The
        privacy manifest and App Store privacy label must match what you send.
  \item \textbf{Transport}: buffer $\to$ persisted queue $\to$ batch flush by
        count, time or app-backgrounding; gzip; retry with backoff; cap the queue
        (drop oldest, \emph{count} drops); \ct{event\_id} makes retries idempotent.
\end{itemize}

\columnbreak

\section{Example — flags, bucketing, exposure}
\begin{lstlisting}[language=SwiftSheet]
import CryptoKit
struct Flags: Codable { var newCheckout = false }
final class FlagStore {              // read once per session
  private(set) var current: Flags
  init(cached: Data?) {              // disk cache, else defaults
    current = cached.flatMap { try? JSONDecoder()
      .decode(Flags.self, from: $0) } ?? Flags() }
  func refresh() async { /* fetch; save for NEXT launch */ }
}
func bucket(_ exp: String, _ user: String) -> Int { // 0..9999
  let d = SHA256.hash(data: Data("\(exp):\(user)".utf8))
  return Int(d.prefix(4).reduce(UInt32(0)) { $0 << 8 | UInt32($1) }
             % 10_000)               // NOT hashValue
}
let arm = bucket("checkout_v2", uid) < 5_000 ? "control" : "test"
analytics.track("experiment_exposure",   // at the screen
                ["experiment": "checkout_v2", "arm": arm])
\end{lstlisting}

\section{Mobile observability}
\begin{itemize}
  \item \textbf{Crash-free users} = users with no crash / active users;
        \textbf{crash-free sessions} = sessions with no crash / sessions. Users
        hides frequency, sessions hides breadth — watch both, per version.
  \item Symbolication: upload dSYMs per build from CI, or stacks stay hex.
        \textbf{Hangs} (iOS hang rate, s/hour) \textasciitilde{} Android
        \textbf{ANR}. \textbf{Traces}: spans for cold start, screen load, API
        calls, trace id propagated to the backend. \textbf{Breadcrumbs}: last N
        screens, requests, logs attached to each crash; log non-fatal errors too.
        (Crashlytics, Sentry, Datadog: examples.)
  \item \textbf{Release gate}: compare the new version with the previous one
        over the same days and a minimum sample; pause the phased release on a
        regression; separate \emph{binary} rollout from \emph{feature} rollout.
\end{itemize}

\section{Traps}
\begin{itemize}
  \trap{Bucketing with Swift \ct{hashValue} — seeded per process: arms change every launch.}
  \trap{Logging exposure at launch for every user — dilutes the effect.}
  \trap{``Phased release lets me roll back'' — it only pauses; ship a fix or flip a flag.}
  \trap{Stopping when p dips under 0.05 (peeking); ignoring SRM.}
  \trap{Randomising a signed-in feature by device: one user sees both arms
        on iPhone and iPad — pick the unit you analyse.}
\end{itemize}

\section{Remember}
\textbf{Default off, hash to bucket, log exposure, guard the rails, ramp in
steps, delete the flag.}


\end{multicols}

\noindent{\footnotesize\color{sheetGrey}\textit{Related:} signing-and-cicd (TestFlight,
phased release) · crashes-symbolication (dSYMs) · logging-signposts-power
(MetricKit, Organizer) · ab-testing-statistics · progressive-delivery · att-adattributionkit}

\end{document}
